\section{Test e Scenari di Bernstein}

In questa sezione vengono descritti i test implementati per il modulo \texttt{bernstein}, volti a garantire la correttezza matematica delle basi di Bernstein e delle utility di calcolo.

\subsection{StandardBernsteinBasisTest}
I test per le basi di Bernstein standard verificano le proprietà fondamentali dei polinomi di Bernstein in un intervallo $[a, b]$:
\begin{itemize}
    \item \textbf{Valori Noti (\texttt{knownValues}):} Verifica che per $n=1$ e $n=2$ i valori calcolati in punti specifici (come gli estremi $0$ e $1$ e il punto medio $0.5$) corrispondano ai valori teorici attesi.
    \item \textbf{Partizione dell'Unità (\texttt{partitionOfUnity}):} Conferma che la somma di tutte le funzioni di base in un punto qualsiasi dell'intervallo sia esattamente $1$ e che ogni singolo valore sia non negativo.
    \item \textbf{Intervallo Invalido (\texttt{invalidInterval}):} Assicura che venga sollevata un'eccezione (\texttt{IllegalArgumentException}) se l'estremo superiore dell'intervallo è minore di quello inferiore.
\end{itemize}

\subsection{ExponentialBernsteinBasisTest}
Questi test riguardano le basi di Bernstein esponenziali, utilizzate per approssimazioni su domini semi-infiniti $[0, \infty)$:
\begin{itemize}
    \item \textbf{Valori Noti (\texttt{knownValues}):} Verifica il comportamento ai limiti: a $x=0$ solo la prima base ($i=0$) deve essere $1$, mentre per $x$ molto grandi l'ultima base ($i=n$) deve tendere a $1$.
    \item \textbf{Partizione dell'Unità (\texttt{partitionOfUnity}):} Valida la proprietà di partizione dell'unità e la non negatività su vari punti del dominio $[0, 20]$.
    \item \textbf{Lambda Invalido (\texttt{invalidLambda}):} Verifica che il parametro di decadimento $\lambda$ sia strettamente positivo.
\end{itemize}

\subsection{MyMathTest}
Vengono testate le funzioni matematiche di supporto, in particolare il calcolo dei coefficienti binomiali:
\begin{itemize}
    \item \textbf{Valori Piccoli (\texttt{binomialCoefficient\_smallValues}):} Correttezza del calcolo per piccoli valori di $n$ e $k$.
    \item \textbf{Grandi Valori (\texttt{binomialCoefficient\_largeValues}):} Utilizzo di \texttt{BigInteger} per gestire coefficienti che superano il limite dei tipi primitivi (es. $\binom{100}{50}$).
    \item \textbf{Casi Limite (\texttt{binomialCoefficient\_edgeCases}):} Gestione dei casi in cui $k > n$ o $k < 0$, che devono restituire $0$.
\end{itemize}


\subsection{Validazione e Test calcol polinomi}
Per garantire la correttezza dell'implementazione dell'algoritmo di de Casteljau, sono stati definiti i seguenti test unitari:

\begin{itemize}
    \item \textbf{Costante:} Verifica che un polinomio con tutti i coefficienti uguali a $C$ valga $C$ in tutto il dominio. Garantisce la proprietà di "partition of unity" della base.
    \item \textbf{Identità:} Utilizzando i coefficienti $\beta_i = \frac{i}{n}$, il polinomio deve coincidere con la funzione identità $f(x)=x$ (linearità).
    \item \textbf{Boundaries:} Verifica la valutazione agli estremi del supporto (es. $x=0$ e $x=1$ per basi lineari). Garantisce che $P(0) = \beta_0$ e $P(1) = \beta_n$.
    \item \textbf{Simmetria:} Verifica che invertendo l'ordine dei coefficienti e valutando in $1-x$, il risultato sia identico.
\end{itemize}
